The Value Proposition explains how the sponsors will benefit from your work, and why they should invest funding, time, and expertise in supporting your team. Here, you are essentially making a case for the project. There are many ways in which value can be returned to your stakeholders (industrial sponsors, instructors, the university, etc.), list any that may help you convince them to "buy in".

\begin{itemize}

\item \textbf{Reusable}: This could be a reusable platform, the SIEM can be reused or extended by future engineers or companies that would like to add features in the future.

\item \textbf{Research Purposes}: It provides a good foundation for further work in cybersecurity analytics.

\item \textbf{Collaboration Potential}: This project allows for collaboration potential in the future, by laying the groundwork for partnerships with different security firms or agencies that would like an insight in modern SIEM research, and new implementations.

\item \textbf{Hands-on Security Platform}: The Sponsor will benefit from us working on SIEM prototype that will integrate open-source tools (Suricata,Zeek), event log pipelines ( Beats, Elasticsearch/Fluentd), and the AI portion of AI-driven analysis 
(Gemini/ChatGPT/DeepSeek R1). Which allow us to be cost effective but still provide real security monitoring without proprietary SIEMs.

\item \textbf{Teaching Aid}: PSIEM can serve as a reusable tool to teach future students in class, giving the instructor a hands-on platform to demonstrate methods that were talked about in class. 

\end{itemize}